\documentclass[11pt,a4paper]{article}
\usepackage[utf8]{inputenc}
\usepackage[T1]{fontenc}
\usepackage[catalan]{babel}
\usepackage{lmodern}
\usepackage[sfdefault,lining]{FiraSans}
\usepackage[a4paper,top=2.3cm,bottom=2.2cm,left=1.8cm,right=1.8cm,headheight=1.3cm,headsep=0.4cm,footskip=1cm]{geometry}
\usepackage{amsmath,amssymb}
\allowdisplaybreaks
\usepackage{icomma}
\usepackage{graphicx}
\usepackage[table]{xcolor}
\usepackage{tikz}
\usepackage{fontawesome5}
\usepackage{ccicons}
\usepackage{fancyhdr}
\usepackage{lastpage}
\usepackage{tasks}
\usepackage{needspace}
\usepackage[most]{tcolorbox}
\usepackage{array}
\usepackage[hidelinks,pdftitle={Fitxa d'exercicis de la unitat 4: Problemes mètrics. 2n de Batxillerat MII},pdfsubject={Matemàtiques II},pdfkeywords={angles, distàncies, perpendicular comuna, projeccions, simètrics, àrees, volums, pla mediador, plans bisectors, esfera, PAU}]{hyperref}

\definecolor{navy}{HTML}{1B2748}
\definecolor{accent}{HTML}{E8603A}
\definecolor{teal}{HTML}{1F9E90}
\definecolor{soft}{HTML}{EEF1F8}
\definecolor{muted}{HTML}{5D6685}

\setlength{\parindent}{0pt}
\setlength{\parskip}{3pt}
\settasks{label=\alph*), label-format=\bfseries\color{navy}, label-width=1.3em, item-indent=1.8em, column-sep=0.6em, after-item-skip=5pt, before-skip=3pt, after-skip=4pt}

% ---------- capçalera i peu ----------
\newcommand{\llicencia}{\href{https://creativecommons.org/licenses/by-nc-sa/4.0/deed.ca}{\ccbyncsa}\ \ Llicència Creative Commons Reconeixement-NoComercial-CompartirIgual 4.0 Internacional (CC BY-NC-SA 4.0)}
\pagestyle{fancy}
\fancyhf{}
\fancyhead[L]{\raisebox{-0.35\height}{\includegraphics[height=0.95cm]{logo-gva.png}}}
\fancyhead[C]{\small\color{muted}Fitxa d'exercicis · Unitat 4: Problemes mètrics · 2n Batxillerat MII}
\fancyhead[R]{\raisebox{-0.35\height}{\includegraphics[height=0.95cm]{logo-vermellar.png}}}
\fancyfoot[L]{\scriptsize\color{muted}\llicencia}
\fancyfoot[R]{\small\color{muted}\thepage}
\renewcommand{\headrulewidth}{0.4pt}
\renewcommand{\headrule}{\color{muted!40}\hrule height \headrulewidth}
\fancypagestyle{primera}{\fancyhf{}\renewcommand{\headrulewidth}{0pt}
  \fancyfoot[L]{\scriptsize\color{muted}\llicencia}
  \fancyfoot[R]{\small\color{muted}\thepage}}

% ---------- nivells i exercicis ----------
\newcommand{\estrelles}[1]{{\color{accent}\ifcase#1\or\faStar\faStar[regular]\faStar[regular]\or\faStar\faStar\faStar[regular]\or\faStar\faStar\faStar\fi}}
\newcounter{exer}
\newcommand{\ex}[1]{\par\needspace{4\baselineskip}\medskip\refstepcounter{exer}\noindent{\color{navy}\bfseries\theexer.}\ \ }

\newcommand{\seccio}[2]{\par\needspace{10\baselineskip}\bigskip
  \begin{tcolorbox}[colback=navy,colframe=navy,coltext=white,arc=2mm,boxrule=0pt,left=3mm,right=3mm,top=1.2mm,bottom=1.2mm]
  {\large\bfseries #1\quad #2}
  \end{tcolorbox}\vspace{-2pt}}


\setlength{\emergencystretch}{3em}  % les fórmules llargues (equacions de rectes, ternes) passen a la línia següent

% quadre de recordatori (només a la fitxa d'exercicis)
\newcommand{\recorda}[1]{\par\smallskip\begin{tcolorbox}[colback=accent!7,colframe=accent!55,arc=2mm,boxrule=0.5pt,left=2.5mm,right=2.5mm,top=1mm,bottom=1mm]\small{\color{accent}\faLightbulb[regular]}\ \textbf{Recorda:}\ #1\end{tcolorbox}}
% etiqueta dels exercicis de la PAU d'altres comunitats
\newcommand{\pau}[1]{{\color{teal}\small\bfseries\faUniversity\ PAU · #1.}\ }
\newcommand{\res}[1]{#1}\newcommand{\rest}[1]{#1}

% =====================================================================
% Capçalera FentMates de les fitxes: un color per a cada curs.
% Els colors són els de la pàgina inicial (LEVELS a assets/js/materials.js);
% el logo ve de _eines/logo/pdf/ (python3 _eines/logo/genera.py).
%
% Ús, al final del preàmbul de la unitat:
%   % =====================================================================
% Capçalera FentMates de les fitxes: un color per a cada curs.
% Els colors són els de la pàgina inicial (LEVELS a assets/js/materials.js);
% el logo ve de _eines/logo/pdf/ (python3 _eines/logo/genera.py).
%
% Ús, al final del preàmbul de la unitat:
%   % =====================================================================
% Capçalera FentMates de les fitxes: un color per a cada curs.
% Els colors són els de la pàgina inicial (LEVELS a assets/js/materials.js);
% el logo ve de _eines/logo/pdf/ (python3 _eines/logo/genera.py).
%
% Ús, al final del preàmbul de la unitat:
%   \input{../comu/fentmates.tex}
%   \fmcurs{1eso}
% i, al començament del document:
%   \capcalera{Matemàtiques}{1r ESO}{Fitxa d'exercicis · Unitat 2}{Divisibilitat}
% Opcional abans de \input: \def\fmneed{14} (espai mínim abans de cada secció).
% =====================================================================
\makeatletter
\def\fm@logodir{../../logo/pdf}
\providecommand{\fmneed}{10}
\def\fm@cursnom{}

% clau del curs, color i color del text damunt del color
\newcommand{\fm@curs}[3]{\@namedef{fm@color@#1}{#2}\@namedef{fm@text@#1}{#3}}
\fm@curs{1eso}{E6B84A}{navy}
\fm@curs{2eso}{3B82F6}{white}
\fm@curs{3eso}{FF7A50}{white}
\fm@curs{4eso-a}{A78BFA}{white}
\fm@curs{4eso-b}{6D28D9}{white}
\fm@curs{1batx-mi}{9B1C1C}{white}
\fm@curs{1batx-ccss}{EF6B6B}{white}
\fm@curs{2batx-mii}{1E6B47}{white}
\fm@curs{2batx-ccss}{2BB3A3}{white}

\newcommand{\fmcurs}[1]{%
  \@ifundefined{fm@color@#1}{\PackageError{fentmates}{Curs desconegut: #1}{}}{}%
  \edef\fm@tmp{\noexpand\definecolor{curs}{HTML}{\@nameuse{fm@color@#1}}%
               \noexpand\colorlet{curstext}{\@nameuse{fm@text@#1}}}\fm@tmp
  \colorlet{cursfosc}{curs!72!black}%
  \def\fm@logo{\fm@logodir/vertical-#1.pdf}%
  \def\fm@marca{\fm@logodir/marca-#1.pdf}}

% ---------- capçalera de la primera pàgina ----------
% #1 matèria (petit), #2 curs, #3 tipus de document i unitat, #4 títol
\newcommand{\capcalera}[4]{%
  \gdef\fm@cursnom{#2}%
  \par\noindent
  \begin{tcolorbox}[colback=curs,colframe=curs,boxrule=0pt,arc=2.2mm,width=3.5cm,height=2.15cm,
      valign=center,left=3mm,right=2mm,top=1mm,bottom=1mm,nobeforeafter,box align=center]
    \color{curstext}{\footnotesize #1}\par\vspace{1pt}{\fontsize{19}{21}\selectfont\bfseries #2\par}%
  \end{tcolorbox}\hspace{5mm}%
  \begin{minipage}[c]{\dimexpr\linewidth-3.5cm-5mm-2.9cm\relax}
    {\color{cursfosc}\small\bfseries\MakeUppercase{#3}}\par\vspace{2pt}
    {\color{navy}\fontsize{20}{23}\selectfont\bfseries #4\par}
  \end{minipage}\hfill
  \begin{minipage}[c]{2.6cm}\raggedleft\includegraphics[width=2.6cm]{\fm@logo}\end{minipage}\par
  \vspace{2pt}}

% ---------- barres de secció del color del curs ----------
\renewcommand{\seccio}[2]{\par\needspace{\fmneed\baselineskip}\bigskip
  \begin{tcolorbox}[colback=curs,colframe=curs,coltext=curstext,arc=2mm,boxrule=0pt,left=3mm,right=3mm,top=1.2mm,bottom=1.2mm]
  {\large\bfseries #1\quad #2}
  \end{tcolorbox}\vspace{-2pt}}

% ---------- estrelles de dificultat del color del curs (a Batxillerat no se'n posen) ----------
\@ifundefined{estrelles}{}{%
  \renewcommand{\estrelles}[1]{{\color{curs}\ifcase#1\or\faStar\faStar[regular]\faStar[regular]\or\faStar\faStar\faStar[regular]\or\faStar\faStar\faStar\fi}}}

% ---------- capçalera i peu de les pàgines ----------
% primera pàgina: logos institucionals al peu
\fancypagestyle{primera}{\fancyhf{}\renewcommand{\headrulewidth}{0pt}
  \fancyfoot[L]{\raisebox{-0.32\height}{\includegraphics[height=0.78cm]{logo-gva.png}}\hspace{5mm}%
    \raisebox{-0.32\height}{\includegraphics[height=0.78cm]{logo-vermellar.png}}}
  \fancyfoot[R]{\scriptsize\color{muted}\href{https://creativecommons.org/licenses/by-nc-sa/4.0/deed.ca}{\ccbyncsa}\ \ CC BY-NC-SA 4.0 · fentmates.org\qquad\small\thepage}}
% la resta: el símbol a l'esquerra i el curs a la dreta, amb el títol al centre com abans
\fancyhead[L]{\raisebox{-0.3\height}{\includegraphics[height=0.62cm]{\fm@marca}}}
\fancyhead[R]{\small\bfseries\color{cursfosc}\fm@cursnom}
\renewcommand{\headrule}{\color{curs!45}\hrule height \headrulewidth}
\makeatother

%   \fmcurs{1eso}
% i, al començament del document:
%   \capcalera{Matemàtiques}{1r ESO}{Fitxa d'exercicis · Unitat 2}{Divisibilitat}
% Opcional abans de \input: \def\fmneed{14} (espai mínim abans de cada secció).
% =====================================================================
\makeatletter
\def\fm@logodir{../../logo/pdf}
\providecommand{\fmneed}{10}
\def\fm@cursnom{}

% clau del curs, color i color del text damunt del color
\newcommand{\fm@curs}[3]{\@namedef{fm@color@#1}{#2}\@namedef{fm@text@#1}{#3}}
\fm@curs{1eso}{E6B84A}{navy}
\fm@curs{2eso}{3B82F6}{white}
\fm@curs{3eso}{FF7A50}{white}
\fm@curs{4eso-a}{A78BFA}{white}
\fm@curs{4eso-b}{6D28D9}{white}
\fm@curs{1batx-mi}{9B1C1C}{white}
\fm@curs{1batx-ccss}{EF6B6B}{white}
\fm@curs{2batx-mii}{1E6B47}{white}
\fm@curs{2batx-ccss}{2BB3A3}{white}

\newcommand{\fmcurs}[1]{%
  \@ifundefined{fm@color@#1}{\PackageError{fentmates}{Curs desconegut: #1}{}}{}%
  \edef\fm@tmp{\noexpand\definecolor{curs}{HTML}{\@nameuse{fm@color@#1}}%
               \noexpand\colorlet{curstext}{\@nameuse{fm@text@#1}}}\fm@tmp
  \colorlet{cursfosc}{curs!72!black}%
  \def\fm@logo{\fm@logodir/vertical-#1.pdf}%
  \def\fm@marca{\fm@logodir/marca-#1.pdf}}

% ---------- capçalera de la primera pàgina ----------
% #1 matèria (petit), #2 curs, #3 tipus de document i unitat, #4 títol
\newcommand{\capcalera}[4]{%
  \gdef\fm@cursnom{#2}%
  \par\noindent
  \begin{tcolorbox}[colback=curs,colframe=curs,boxrule=0pt,arc=2.2mm,width=3.5cm,height=2.15cm,
      valign=center,left=3mm,right=2mm,top=1mm,bottom=1mm,nobeforeafter,box align=center]
    \color{curstext}{\footnotesize #1}\par\vspace{1pt}{\fontsize{19}{21}\selectfont\bfseries #2\par}%
  \end{tcolorbox}\hspace{5mm}%
  \begin{minipage}[c]{\dimexpr\linewidth-3.5cm-5mm-2.9cm\relax}
    {\color{cursfosc}\small\bfseries\MakeUppercase{#3}}\par\vspace{2pt}
    {\color{navy}\fontsize{20}{23}\selectfont\bfseries #4\par}
  \end{minipage}\hfill
  \begin{minipage}[c]{2.6cm}\raggedleft\includegraphics[width=2.6cm]{\fm@logo}\end{minipage}\par
  \vspace{2pt}}

% ---------- barres de secció del color del curs ----------
\renewcommand{\seccio}[2]{\par\needspace{\fmneed\baselineskip}\bigskip
  \begin{tcolorbox}[colback=curs,colframe=curs,coltext=curstext,arc=2mm,boxrule=0pt,left=3mm,right=3mm,top=1.2mm,bottom=1.2mm]
  {\large\bfseries #1\quad #2}
  \end{tcolorbox}\vspace{-2pt}}

% ---------- estrelles de dificultat del color del curs (a Batxillerat no se'n posen) ----------
\@ifundefined{estrelles}{}{%
  \renewcommand{\estrelles}[1]{{\color{curs}\ifcase#1\or\faStar\faStar[regular]\faStar[regular]\or\faStar\faStar\faStar[regular]\or\faStar\faStar\faStar\fi}}}

% ---------- capçalera i peu de les pàgines ----------
% primera pàgina: logos institucionals al peu
\fancypagestyle{primera}{\fancyhf{}\renewcommand{\headrulewidth}{0pt}
  \fancyfoot[L]{\raisebox{-0.32\height}{\includegraphics[height=0.78cm]{logo-gva.png}}\hspace{5mm}%
    \raisebox{-0.32\height}{\includegraphics[height=0.78cm]{logo-vermellar.png}}}
  \fancyfoot[R]{\scriptsize\color{muted}\href{https://creativecommons.org/licenses/by-nc-sa/4.0/deed.ca}{\ccbyncsa}\ \ CC BY-NC-SA 4.0 · fentmates.org\qquad\small\thepage}}
% la resta: el símbol a l'esquerra i el curs a la dreta, amb el títol al centre com abans
\fancyhead[L]{\raisebox{-0.3\height}{\includegraphics[height=0.62cm]{\fm@marca}}}
\fancyhead[R]{\small\bfseries\color{cursfosc}\fm@cursnom}
\renewcommand{\headrule}{\color{curs!45}\hrule height \headrulewidth}
\makeatother

%   \fmcurs{1eso}
% i, al començament del document:
%   \capcalera{Matemàtiques}{1r ESO}{Fitxa d'exercicis · Unitat 2}{Divisibilitat}
% Opcional abans de \input: \def\fmneed{14} (espai mínim abans de cada secció).
% =====================================================================
\makeatletter
\def\fm@logodir{../../logo/pdf}
\providecommand{\fmneed}{10}
\def\fm@cursnom{}

% clau del curs, color i color del text damunt del color
\newcommand{\fm@curs}[3]{\@namedef{fm@color@#1}{#2}\@namedef{fm@text@#1}{#3}}
\fm@curs{1eso}{E6B84A}{navy}
\fm@curs{2eso}{3B82F6}{white}
\fm@curs{3eso}{FF7A50}{white}
\fm@curs{4eso-a}{A78BFA}{white}
\fm@curs{4eso-b}{6D28D9}{white}
\fm@curs{1batx-mi}{9B1C1C}{white}
\fm@curs{1batx-ccss}{EF6B6B}{white}
\fm@curs{2batx-mii}{1E6B47}{white}
\fm@curs{2batx-ccss}{2BB3A3}{white}

\newcommand{\fmcurs}[1]{%
  \@ifundefined{fm@color@#1}{\PackageError{fentmates}{Curs desconegut: #1}{}}{}%
  \edef\fm@tmp{\noexpand\definecolor{curs}{HTML}{\@nameuse{fm@color@#1}}%
               \noexpand\colorlet{curstext}{\@nameuse{fm@text@#1}}}\fm@tmp
  \colorlet{cursfosc}{curs!72!black}%
  \def\fm@logo{\fm@logodir/vertical-#1.pdf}%
  \def\fm@marca{\fm@logodir/marca-#1.pdf}}

% ---------- capçalera de la primera pàgina ----------
% #1 matèria (petit), #2 curs, #3 tipus de document i unitat, #4 títol
\newcommand{\capcalera}[4]{%
  \gdef\fm@cursnom{#2}%
  \par\noindent
  \begin{tcolorbox}[colback=curs,colframe=curs,boxrule=0pt,arc=2.2mm,width=3.5cm,height=2.15cm,
      valign=center,left=3mm,right=2mm,top=1mm,bottom=1mm,nobeforeafter,box align=center]
    \color{curstext}{\footnotesize #1}\par\vspace{1pt}{\fontsize{19}{21}\selectfont\bfseries #2\par}%
  \end{tcolorbox}\hspace{5mm}%
  \begin{minipage}[c]{\dimexpr\linewidth-3.5cm-5mm-2.9cm\relax}
    {\color{cursfosc}\small\bfseries\MakeUppercase{#3}}\par\vspace{2pt}
    {\color{navy}\fontsize{20}{23}\selectfont\bfseries #4\par}
  \end{minipage}\hfill
  \begin{minipage}[c]{2.6cm}\raggedleft\includegraphics[width=2.6cm]{\fm@logo}\end{minipage}\par
  \vspace{2pt}}

% ---------- barres de secció del color del curs ----------
\renewcommand{\seccio}[2]{\par\needspace{\fmneed\baselineskip}\bigskip
  \begin{tcolorbox}[colback=curs,colframe=curs,coltext=curstext,arc=2mm,boxrule=0pt,left=3mm,right=3mm,top=1.2mm,bottom=1.2mm]
  {\large\bfseries #1\quad #2}
  \end{tcolorbox}\vspace{-2pt}}

% ---------- estrelles de dificultat del color del curs (a Batxillerat no se'n posen) ----------
\@ifundefined{estrelles}{}{%
  \renewcommand{\estrelles}[1]{{\color{curs}\ifcase#1\or\faStar\faStar[regular]\faStar[regular]\or\faStar\faStar\faStar[regular]\or\faStar\faStar\faStar\fi}}}

% ---------- capçalera i peu de les pàgines ----------
% primera pàgina: logos institucionals al peu
\fancypagestyle{primera}{\fancyhf{}\renewcommand{\headrulewidth}{0pt}
  \fancyfoot[L]{\raisebox{-0.32\height}{\includegraphics[height=0.78cm]{logo-gva.png}}\hspace{5mm}%
    \raisebox{-0.32\height}{\includegraphics[height=0.78cm]{logo-vermellar.png}}}
  \fancyfoot[R]{\scriptsize\color{muted}\href{https://creativecommons.org/licenses/by-nc-sa/4.0/deed.ca}{\ccbyncsa}\ \ CC BY-NC-SA 4.0 · fentmates.org\qquad\small\thepage}}
% la resta: el símbol a l'esquerra i el curs a la dreta, amb el títol al centre com abans
\fancyhead[L]{\raisebox{-0.3\height}{\includegraphics[height=0.62cm]{\fm@marca}}}
\fancyhead[R]{\small\bfseries\color{cursfosc}\fm@cursnom}
\renewcommand{\headrule}{\color{curs!45}\hrule height \headrulewidth}
\makeatother

\fmcurs{2batx-mii}

\begin{document}
\thispagestyle{primera}

\capcalera{Matemàtiques II}{2n Batx}{Fitxa d'exercicis · Unitat 4}{\raggedright Problemes mètrics}
\vspace{10pt}
Nom i cognoms: \rule{8.4cm}{0.4pt}\hfill Grup: \rule{1.6cm}{0.4pt}\hfill Data: \rule{2.2cm}{0.4pt}
\vspace{8pt}

\begin{tcolorbox}[colback=soft,colframe=soft,arc=2mm,left=3mm,right=3mm,top=1.5mm,bottom=1.5mm]
\small Resol els exercicis al quadern justificant cada pas. Fes un esbós quan t'ajude a veure la situació i comprova sempre els resultats (que la projecció és del pla o de la recta, que el punt mig del segment $PP'$ és el punt $M$\ldots). Els exercicis marcats amb {\color{teal}\bfseries\faUniversity\ PAU} són de les proves d'accés a la universitat d'altres comunitats autònomes (2024--2026), traduïts al valencià.
\end{tcolorbox}

% =====================================================
\seccio{1}{Angles entre rectes i plans}

\ex{1} Calcula l'angle que formen les rectes:\begin{tasks}(1)
  \task \mbox{$r:\dfrac{x-1}{2}=\dfrac{y+1}{1}=\dfrac{z}{-2}$} \ i \ $s:(x,\,y,\,z)=(0,\,1,\,3)+\lambda\,(4,\,0,\,3)$
  \task $r:\left\{\begin{aligned}x+y&=1\\y-z&=2\end{aligned}\right.$ \ i \ $s:x=y=z$
\end{tasks}\recorda{\mbox{$\cos\alpha=\dfrac{|\vec d_r\cdot\vec d_s|}{|\vec d_r|\,|\vec d_s|}$}, sempre entre $0^\circ$ i $90^\circ$.}

\ex{1} Calcula l'angle que forma la recta \mbox{$r:\dfrac{x-2}{1}=\dfrac{y}{-1}=\dfrac{z+1}{2}$} amb cadascun dels plans següents:\begin{tasks}(2)
  \task $\pi_1:x+y+z=4$
  \task $\pi_2:2x-2y+4z=1$
\end{tasks}

\ex{1} Calcula l'angle que formen els plans:\begin{tasks}(2)
  \task $x+y=3$ \ i \ $x+z=1$
  \task $2x-y+2z=0$ \ i el pla $OXY$
\end{tasks}

\ex{2} Calcula els angles que forma la recta de vector director $(2,\,-1,\,2)$ amb cadascun dels eixos de coordenades.

\ex{2} Troba $m$ perquè la recta $(x,\,y,\,z)=(1,\,0,\,0)+\lambda\,(1,\,m,\,1)$ forme un angle de $30^\circ$ amb el pla $z=0$.

\ex{2} Hi ha algun valor de $k$ per al qual els plans $x+ky+z=0$ i $2x-y+kz=3$ siguen perpendiculars? I paral·lels?

\ex{2} \pau{Canàries, extraordinària 2026} Se sap que el pla $\pi:x+3y+2z=10$ i la recta \mbox{$r:\dfrac{x+21}{11}=3-y=\dfrac{5-z}{2}$} són secants.\begin{tasks}(1)
  \task Troba el punt de tall del pla $\pi$ i la recta $r$.
  \task Calcula l'angle que formen el pla $\pi$ i la recta $r$.
\end{tasks}

\ex{2} \pau{Castella-la Manxa, ordinària 2024} Siguen els vectors $\vec u=(1,\,a,\,a)$ i $\vec v=(-1,\,0,\,2)$, amb $a\in\mathbb R$. Determina el valor de $a$ perquè l'angle entre els vectors $\vec u$ i $\vec v$ siga de $60^\circ$.

\ex{2} \pau{País Basc, extraordinària 2025} Siguen $\pi$ el pla que passa pels punts $A(2,\,3,\,4)$, $B(3,\,-1,\,-2)$, $C(5,\,-1,\,2)$ i $r$ la recta que passa pels punts $D(6,\,-5,\,-4)$, $E(7,\,1,\,4)$.\begin{tasks}(1)
  \task Calcula l'angle entre el pla $\pi$ i la recta $r$, expressant el resultat en graus, minuts i segons.
  \task En cas que $r$ i $\pi$ es tallen, calcula el punt d'intersecció. En cas contrari, calcula la distància entre la recta $r$ i el pla $\pi$.
\end{tasks}

\ex{2} \pau{Canàries, extraordinària 2024}\begin{tasks}(1)
  \task Justifica si poden existir vectors $\vec u$ i $\vec v$, amb el mateix origen, que complisquen $|\vec u|=2$, $|\vec v|=3$ i $\vec u\cdot\vec v=8$.
  \task Donats el pla $\pi:x+3y+2z+3=0$ i la recta $r:\left\{\begin{aligned}2x-3y-z&=4\\x+y+2z&=-3\end{aligned}\right.$, que són secants, calcula el punt de tall de la recta i el pla i l'angle que formen.
\end{tasks}

\ex{2} \pau{La Rioja, ordinària 2024} Troba l'equació contínua de la recta $s$ que està continguda en el pla $\pi:x+y-2z+1=0$ i talla perpendicularment la recta $r:\left\{\begin{aligned}x+y+z&=-1\\4x-y+z&=3\end{aligned}\right.$


% =====================================================
\seccio{2}{Distància entre punts i d'un punt a un pla}

\ex{1} Calcula:\begin{tasks}(1)
  \task la distància del punt $P(2,\,-1,\,3)$ al pla $3x-6y+2z+1=0$;
  \task la distància de l'origen al pla \mbox{$\dfrac x2+\dfrac y3+\dfrac z6=1$};
  \task la distància entre els punts $A(1,\,2,\,-1)$ i $B(3,\,-1,\,5)$.
\end{tasks}\recorda{\mbox{$d(P,\pi)=\dfrac{|Ax_0+By_0+Cz_0+D|}{\sqrt{A^2+B^2+C^2}}$}}

\ex{1} Troba la projecció ortogonal del punt $P(3,\,3,\,-1)$ sobre el pla $\pi:x+y-z=1$ i comprova que la distància de $P$ a la projecció és $d(P,\pi)$.

\ex{2} Troba els punts de l'eix $OZ$ que disten 3 unitats del pla $2x+y-2z=1$.

\ex{2} Troba $a$ perquè el punt $P(1,\,a,\,1)$ diste 2 unitats del pla $x+2y+2z-3=0$.

\ex{2} \pau{Astúries, ordinària 2026} Es considera el pla $\pi:x+2y+3z=6$ i el punt $P(0,\,-1,\,-2)$.\begin{tasks}(1)
  \task Troba l'equació d'una recta paral·lela a $\pi$ que continga el punt $P$. És única aquesta recta?
  \task Calcula el punt $Q\in\pi$ que dista menys de $P$. Quina distància hi ha entre $P$ i $Q$?
\end{tasks}

\ex{2} \pau{Catalunya, extraordinària 2024} Considera el punt $P(1,\,3,\,0)$ i el pla $\pi:x+2y-2z=-7$.\begin{tasks}(1)
  \task Siga $r$ la recta perpendicular a $\pi$ que passa per $P$. Calcula el punt d'intersecció de $\pi$ amb $r$.
  \task Calcula la distància $d$ del punt $P$ al pla $\pi$.
  \task Calcula l'equació d'un altre pla $\pi'$ paral·lel a $\pi$ que també estiga a distància $d$ de $P$.
\end{tasks}

\ex{2} \pau{Galícia, ordinària 2026} Considera els punts $P(1,\,0,\,1)$, $Q(0,\,1,\,0)$ i $R(1,\,1,\,1)$.\begin{tasks}(1)
  \task Demostra que $P$, $Q$ i $R$ no estan alineats i obtín l'equació implícita del pla $\pi$ que passa per ells.
  \task Calcula el valor o valors de $a$ que fan que la distància del punt $S(a,\,0,\,0)$ a aquest pla siga igual a $\sqrt2$ unitats de longitud.
\end{tasks}

\ex{2} \pau{Astúries, extraordinària 2026} Donat el pla $\pi:2x+2y+z=5$:\begin{tasks}(1)
  \task Troba l'equació d'un pla paral·lel a $\pi$ la distància del qual a l'origen siga 2. Hi ha més d'un pla que complisca aquestes condicions?
  \task Calcula el punt $P$ del pla $\pi$ que està més pròxim a l'origen $O$.
  \task Calcula l'àrea del paral·lelogram format pels segments $OP$ i $OQ$, on $Q(3,\,3,\,3)$.
\end{tasks}

\ex{2} \pau{Andalusia, ordinària 2025} Siguen els punts $O(0,\,0,\,0)$, $A(0,\,2,\,-2)$, $B(1,\,2,\,m)$ i $C(2,\,3,\,2)$.\begin{tasks}(1)
  \task Troba els valors de $m$ perquè el tetraedre determinat pels punts $O$, $A$, $B$ i $C$ tinga un volum de 3 unitats cúbiques.
  \task Per a $m=0$, calcula la distància del punt $O$ al pla que passa pels punts $A$, $B$ i $C$.
\end{tasks}


% =====================================================
\seccio{3}{Distància d'un punt a una recta}

\ex{1} Calcula la distància:\begin{tasks}(1)
  \task del punt $P(1,\,2,\,3)$ a la recta $x=y=z$;
  \task del punt $P(0,\,1,\,-1)$ a la recta $r:\left\{\begin{aligned}x-y&=0\\z&=2\end{aligned}\right.$
\end{tasks}\recorda{\mbox{$d(P,r)=\dfrac{|\overrightarrow{AP}\times\vec d|}{|\vec d|}$}, amb $A\in r$ i $\vec d$ un vector director de $r$.}

\ex{2} Troba el punt $Q$ de la recta $r:(x,\,y,\,z)=(1,\,0,\,1)+\lambda\,(1,\,1,\,1)$ més pròxim a $P(3,\,1,\,4)$, la distància de $P$ a $r$ i la recta que passa per $P$ i talla perpendicularment $r$.

\ex{2} Troba els punts de la recta $r:(x,\,y,\,z)=(1,\,1,\,0)+\lambda\,(2,\,1,\,2)$ que estan a 3 unitats del punt $A(1,\,1,\,0)$.

\ex{2} Calcula l'altura del triangle de vèrtexs $A(1,\,0,\,0)$, $B(0,\,1,\,0)$ i $C(1,\,1,\,1)$ corresponent al costat $AB$.

\ex{2} \pau{La Rioja, extraordinària 2026} Siguen la recta \mbox{$r:\dfrac x3=y=\dfrac{z-11}{-1}$} i el punt $P(0,\,1,\,1)$.\begin{tasks}(1)
  \task Determina l'equació del pla que passa pel punt $P$ i és perpendicular a la recta $r$.
  \task Troba la distància de $P$ a la recta $r$.
\end{tasks}

\ex{2} \pau{Cantàbria, ordinària 2025} Considera la recta $r:\left\{\begin{aligned}2x-y&=3\\y-2z&=1\end{aligned}\right.$ \ i el punt $P(1,\,1,\,1)$.\begin{tasks}(1)
  \task Determina els punts de $r$ que estan a una distància de $\sqrt{14}$ unitats de $P$.
  \task Obtín l'equació del pla que conté $r$ i $P$.
  \task Calcula la distància entre $r$ i $P$.
\end{tasks}

\ex{2} \pau{Extremadura, extraordinària 2025} Donats la recta \mbox{$s:\dfrac{x}{-4}=\dfrac{y-1}{3}=\dfrac{z+1}{0}$}, el pla $\pi:x-2y+3z-6=0$ i el punt $P(1,\,-1,\,0)$:\begin{tasks}(1)
  \task Obtín l'equació del pla perpendicular a la recta $s$ que passa per $P$.
  \task Calcula la distància del punt $P$ a la recta $s$.
  \task Calcula l'angle que forma la recta $s$ amb el pla $\pi$.
\end{tasks}

\ex{2} \pau{Cantàbria, ordinària 2026} Considera els punts $A(0,\,0,\,4)$, $B(-3,\,0,\,0)$ i $C(0,\,-2,\,0)$.\begin{tasks}(1)
  \task Troba l'equació general del pla que conté $A$, $B$ i $C$.
  \task Calcula la distància entre el punt $P(1,\,1,\,1)$ i la recta que conté $A$ i $B$.
  \task Calcula el volum del tetraedre de vèrtexs $A$, $B$, $C$ i $D(0,\,0,\,0)$.
\end{tasks}


% =====================================================
\seccio{4}{Distància entre dues rectes}

\ex{1} Comprova que les rectes $r:(x,\,y,\,z)=(1,\,1,\,0)+\lambda\,(1,\,-1,\,1)$ i $s:(x,\,y,\,z)=(2,\,0,\,3)+\mu\,(-1,\,1,\,-1)$ són paral·leles i calcula la distància entre elles.

\ex{2} \label{ex:pc}Comprova que les rectes $r:\left\{\begin{aligned}x-y&=0\\z&=0\end{aligned}\right.$ \ i \ $s:\left\{\begin{aligned}x&=1\\y-z&=0\end{aligned}\right.$ \ s'encreuen i calcula la distància entre elles.\recorda{Si s'encreuen, \mbox{$d(r,s)=\dfrac{\big|[\overrightarrow{PQ},\vec d_r,\vec d_s]\big|}{|\vec d_r\times\vec d_s|}$}.}

\ex{2} Troba $a$ perquè la distància entre l'eix $OX$ i la recta $s:(x,\,y,\,z)=(0,\,1,\,a)+\mu\,(1,\,1,\,0)$ siga 2.

\ex{2} \pau{País Basc, ordinària 2025} Es consideren la recta $r$ que passa pels punts $A(-4,\,4,\,1)$ i $B(1,\,0,\,-1)$ i la recta $s$ que passa pels punts $C(-2,\,7,\,-1)$ i $D(2,\,3,\,-1)$.\begin{tasks}(1)
  \task Calcula la posició relativa de la recta $r$ respecte de la recta $s$.
  \task En cas que siguen paral·leles o s'encreuen, calcula la distància entre elles. Si es tallen, calcula el punt d'intersecció.
\end{tasks}

\ex{2} \pau{La Rioja, extraordinària 2025} Calcula la distància entre les rectes $r:\left\{\begin{aligned}x&=5+\lambda\\y&=-1\\z&=8+2\lambda\end{aligned}\right.$ \ i \ $s:\left\{\begin{aligned}x&=2+3\mu\\y&=2-\mu\\z&=-1+4\mu\end{aligned}\right.$


% =====================================================
\seccio{5}{Distància entre recta i pla i entre dos plans}

\ex{1} Comprova que la recta $r:x=2+3\lambda,\ y=-\lambda,\ z=1+\lambda$ és paral·lela al pla $\pi:x+2y-z+4=0$ i calcula la distància entre ells.

\ex{1} Calcula la distància entre els plans $3x-4y+12z=5$ i $6x-8y+24z+16=0$.

\ex{2} Troba els plans paral·lels a $2x+y-2z+1=0$ que disten 3 unitats de l'origen.

\ex{2} Troba $m$ perquè la recta $(x,\,y,\,z)=(1,\,1,\,1)+\lambda\,(m,\,1,\,2)$ siga paral·lela al pla $x+2y-2z=0$ i calcula aleshores la distància entre la recta i el pla.

\ex{2} \pau{Andalusia, extraordinària 2025} Considera la recta $r:\left\{\begin{aligned}x-y+z&=3\\x+2y-z&=4\end{aligned}\right.$ \ i el pla $\pi:mx-y-2z=5$.\begin{tasks}(1)
  \task Troba $m$ perquè $r$ i $\pi$ siguen paral·lels.
  \task Per a $m=-8$, calcula la distància de la recta $r$ al pla $\pi$.
\end{tasks}

\ex{2} \pau{Extremadura, ordinària 2025}\begin{tasks}(1)
  \task Comprova que el pla $\pi:x+y-z=3$ i la recta \mbox{$r:\dfrac x3=\dfrac{y-1}{-1}=\dfrac{z+1}{2}$} no es tallen.
  \task Calcula la distància entre el pla $\pi$ i la recta $r$.
  \task Obtín l'equació del pla perpendicular a la recta $r$ que passa pel punt $(0,\,1,\,-1)$.
\end{tasks}

\ex{2} \pau{Galícia, extraordinària 2026} Considera la recta \mbox{$r:\dfrac{x-1}{2}=\dfrac{y-1}{1}=\dfrac z3$}.\begin{tasks}(1)
  \task Calcula l'equació implícita o general del pla $\pi$ que conté la recta $r$ i el punt $Q(2,\,-1,\,1)$.
  \task Comprova que la recta $r$ és paral·lela al pla $\pi':x+y-z-1=0$. Calcula la distància de $r$ a $\pi'$.
\end{tasks}

\ex{2} \pau{Múrcia, ordinària 2026} Considera el pla $\pi:2x-y+3z-5=0$ i la recta \mbox{$r:\dfrac{x-1}{2}=\dfrac{y+1}{1}=\dfrac{z-2}{-1}$}.\begin{tasks}(1)
  \task Estudia la posició relativa de $r$ i $\pi$.
  \task Si $r$ talla $\pi$ en un sol punt, calcula aquest punt i l'angle que forma $r$ amb $\pi$. Si $r$ és paral·lela a $\pi$, calcula la distància de $r$ a $\pi$.
\end{tasks}

\ex{2} \pau{Castella i Lleó, extraordinària 2025} Siguen el pla $\pi:x+y-z=2$ i la recta \mbox{$r:\dfrac{x-1}{-3}=\dfrac y2=\dfrac{z-2}{-1}$}. Calcula l'equació d'un pla $\pi'$ paral·lel al pla $\pi$ que estiga a una distància de $2\sqrt3$ unitats de la recta $r$. És únic aquest pla? Justifica la resposta.

\ex{2} \pau{Catalunya, extraordinària 2025} Considera el pla $\pi:2x-y+z=5$ i el punt $P(0,\,1,\,3)$.\begin{tasks}(1)
  \task Comprova que la distància del punt $P$ al pla $\pi$ és \mbox{$\dfrac{\sqrt6}{2}$}.
  \task Troba l'equació general d'un pla $\pi_1$ paral·lel a $\pi$ que passe pel punt $P$. Quant val la distància entre $\pi_1$ i $\pi$?
  \task Troba l'equació general d'un segon pla $\pi_2$, diferent de $\pi_1$, que siga paral·lel a $\pi$ i que estiga a una distància \mbox{$\dfrac{\sqrt6}{2}$} de $\pi$.
\end{tasks}


% =====================================================
\seccio{6}{Recta que passa per un punt i talla dues rectes}

\ex{1} Troba la recta que passa pel punt $P(2,\,1,\,1)$ i talla les rectes $r:\left\{\begin{aligned}x&=0\\y&=0\end{aligned}\right.$ \ i \ $s:\left\{\begin{aligned}x&=1\\z&=0\end{aligned}\right.$ \ Calcula també els punts de tall.\recorda{$t=\pi_1\cap\pi_2$, on $\pi_1$ és el pla que conté $P$ i $r$, i $\pi_2$, el que conté $P$ i $s$. Comprova que $t$ no és paral·lela a $r$ ni a $s$.}

\ex{2} Troba la recta paral·lela al vector $\vec v=(1,\,0,\,1)$ que talla les rectes $r$ i $s$ de l'exercici anterior.

\ex{2} Troba la recta continguda en el pla $\pi:x+y+z=3$ que passa pel punt $P(1,\,1,\,1)$ i talla la recta $r:(x,\,y,\,z)=\lambda\,(1,\,2,\,0)$.

\ex{2} Raona que no hi ha cap recta que passe per l'origen i talle les rectes $r:(x,\,y,\,z)=(1,\,0,\,0)+\lambda\,(0,\,1,\,0)$ i $s:(x,\,y,\,z)=(0,\,0,\,1)+\mu\,(1,\,0,\,0)$.

\ex{2} \pau{Navarra, extraordinària 2026} Troba l'equació contínua de la recta $t$ que passa pel punt $A(2,\,4,\,-2)$ i talla les rectes \mbox{$s:\dfrac{x-2}{-1}=\dfrac{y+1}{2}=\dfrac z1$} \ i \ $r:\left\{\begin{aligned}x&=-1+\lambda\\y&=-2\\z&=1+3\lambda\end{aligned}\right.$

\ex{2} \pau{Castella i Lleó, extraordinària 2026} Donats el pla $\pi:x+y=0$, el punt $A(1,\,-1,\,0)$ i la recta $s:x=y=z$, troba la recta $r$, en forma paramètrica, que passa per $A$, està continguda en $\pi$ i talla la recta $s$.

\ex{2} \pau{Madrid, extraordinària 2024} Donats el punt $P(5,\,-1,\,2)$ i les rectes \mbox{$r:\dfrac{x-2}{3}=\dfrac{y+1}{-1}=\dfrac{z}{1}$} \ i \ $s:\left\{\begin{aligned}x-y&=5\\x+z&=3\end{aligned}\right.$\begin{tasks}(1)
  \task Estudia la posició relativa de les dues rectes i troba la distància entre elles.
  \task Determina una equació de la recta que passa per $P$ i talla perpendicularment la recta $r$.
\end{tasks}


% =====================================================
\seccio{7}{Perpendicular comuna}

\ex{2} Troba els punts $X\in r$ i $Y\in s$ més pròxims de les rectes $r:\left\{\begin{aligned}x-y&=0\\z&=0\end{aligned}\right.$ \ i \ $s:\left\{\begin{aligned}x&=1\\y-z&=0\end{aligned}\right.$ \ de l'exercici \ref{ex:pc} i la seua perpendicular comuna. Comprova que $|\overrightarrow{XY}|=d(r,s)$.\recorda{Amb els punts genèrics $X\in r$, $Y\in s$, imposa $\overrightarrow{XY}\perp\vec d_r$ i $\overrightarrow{XY}\perp\vec d_s$.}

\ex{2} Siguen $r:(x,\,y,\,z)=(1,\,0,\,0)+\lambda\,(1,\,1,\,0)$ i $s:(x,\,y,\,z)=(0,\,0,\,1)+\mu\,(1,\,-1,\,0)$. Troba la perpendicular comuna pel mètode dels plans i comprova el resultat amb els punts genèrics. Quina és la distància entre $r$ i $s$?\recorda{$\vec w=\vec d_r\times\vec d_s$; \ $\pi_r$ conté $r$ i $\vec w$; \ $\pi_s$ conté $s$ i $\vec w$; \ la perpendicular comuna és $\pi_r\cap\pi_s$.}

\ex{2} \pau{Múrcia, ordinària 2025} Considera les rectes $r:\left\{\begin{aligned}x&=1\\y&=5+\lambda\\z&=\lambda\end{aligned}\right.$ \ i \ $s:\left\{\begin{aligned}x&=7+3\mu\\y&=-5+\mu\\z&=7\end{aligned}\right.$\begin{tasks}(1)
  \task Comprova que s'encreuen.
  \task Troba l'equació de la recta perpendicular a totes dues.
  \task Calcula la distància entre $r$ i $s$.
\end{tasks}

\ex{2} \pau{Navarra, ordinària 2026} Troba l'equació contínua de la recta $t$ que talla perpendicularment les rectes \mbox{$r:\dfrac x1=\dfrac{y+3}{1}=\dfrac{z-2}{-2}$} \ i \ $s:\left\{\begin{aligned}x&=-4+3\lambda\\y&=1+\lambda\\z&=-1+2\lambda\end{aligned}\right.$

\ex{2} \pau{Navarra, extraordinària 2025} Calcula l'equació contínua de la recta $t$ que talla perpendicularment les rectes $r:\left\{\begin{aligned}x+y-2z-10&=0\\x+z-3&=0\end{aligned}\right.$ \ i \ \mbox{$s:\dfrac{x-4}{-1}=\dfrac{y+1}{1}=\dfrac{z}{-1}$}


% =====================================================
\seccio{8}{Simètrics respecte d'un punt i d'una recta}

\ex{1} Calcula el simètric de $A(1,\,-2,\,3)$ respecte del punt $B(0,\,1,\,-1)$. Troba també el punt $M$ respecte del qual $A$ i $A'(5,\,0,\,1)$ són simètrics.

\ex{2} Calcula el simètric del punt $P(2,\,1,\,0)$ respecte de la recta $r:(x,\,y,\,z)=(0,\,0,\,1)+\lambda\,(1,\,0,\,1)$.\recorda{Pla per $P$ perpendicular a $r$ $\to$ $M=r\cap\pi$ $\to$ $P'=2M-P$.}

\ex{2} Calcula el simètric de:\begin{tasks}(1)
  \task $P(1,\,2,\,3)$ respecte de l'eix $OX$;
  \task $P(0,\,0,\,1)$ respecte de la recta $r:\left\{\begin{aligned}x+y&=2\\z&=0\end{aligned}\right.$
\end{tasks}

\ex{2} \pau{Castella i Lleó, ordinària 2024} Troba el punt simètric del punt $P(1,\,0,\,-1)$ respecte de la recta \mbox{$r:\dfrac{x-1}{1}=\dfrac y2=\dfrac z2$}.

\ex{2} \pau{País Basc, extraordinària 2025} Troba el punt $P'$ simètric de $P(4,\,-3,\,0)$ respecte de la recta $r$ que passa pels punts $A(-2,\,0,\,1)$ i $B(-2,\,-1,\,0)$.


% =====================================================
\seccio{9}{Simètrics respecte d'un pla}

\ex{1} Calcula el simètric del punt $P(3,\,1,\,2)$ respecte del pla $\pi:2x-y+z-1=0$.\recorda{Recta per $P$ amb director $\vec n$ $\to$ $M=s\cap\pi$ (la projecció) $\to$ $P'=2M-P$.}

\ex{1} Calcula els simètrics del punt $A(1,\,2,\,3)$ respecte dels plans $z=0$ i $x=y$.

\ex{2} Troba el pla respecte del qual són simètrics els punts $P(1,\,2,\,-1)$ i $P'(3,\,0,\,1)$.

\ex{3} Siguen $r:(x,\,y,\,z)=(1,\,1,\,1)+\lambda\,(1,\,2,\,0)$ i $\pi:x+y+z=0$.\begin{tasks}(1)
  \task Calcula la recta simètrica de $r$ respecte de $\pi$.
  \task Calcula la projecció ortogonal de $r$ sobre $\pi$.
\end{tasks}

\ex{2} Calcula la recta simètrica de $r:(x,\,y,\,z)=(2,\,0,\,0)+\lambda\,(0,\,1,\,-1)$ respecte del pla $y+z=4$.

\ex{2} \pau{Andalusia, ordinària 2024}\begin{tasks}(1)
  \task Troba el punt simètric de $P(2,\,2,\,1)$ respecte de la recta $r:\left\{\begin{aligned}x-2y+z&=2\\y-z&=1\end{aligned}\right.$
  \task Troba el punt simètric de $Q(1,\,-1,\,-3)$ respecte del pla $\pi:x-2y+z+6=0$.
\end{tasks}

\ex{2} \pau{Andalusia, extraordinària 2026} Considera la recta $r:\left\{\begin{aligned}x&=-2\lambda\\y&=-1+2\lambda\\z&=3+2\lambda\end{aligned}\right.$ \ i el pla $\pi:x-2y+z-3=0$.\begin{tasks}(1)
  \task Determina l'equació de la recta que talla perpendicularment $r$ i està continguda en $\pi$.
  \task Troba el punt simètric de $P(1,\,1,\,1)$ respecte de $\pi$.
\end{tasks}

\ex{2} \pau{Canàries, ordinària 2025} Es consideren el punt $A(1,\,0,\,2)$, el pla $\pi:-x+2y+z+1=0$ i la recta $r:\left\{\begin{aligned}4x-7y+2z&=7\\y-2z+5&=0\end{aligned}\right.$\begin{tasks}(1)
  \task Troba la posició relativa del pla $\pi$ i la recta $r$.
  \task Troba el punt simètric de $A$ respecte del pla $\pi$.
\end{tasks}

\ex{2} \pau{Extremadura, extraordinària 2026} Donats els plans $\pi_1:x+ky-3z-6=0$, $\pi_2:2x-y-5=0$ i el punt $A(2,\,3,\,-4)$:\begin{tasks}(1)
  \task Calcula el valor de $k$ perquè els plans $\pi_1$ i $\pi_2$ siguen perpendiculars.
  \task Per a $k=2$, calcula el simètric del punt $A$ respecte del pla $\pi_1$.
  \task Per a $k=2$, calcula la distància del punt $A$ al pla $\pi_1$.
\end{tasks}

\ex{2} \pau{País Basc, extraordinària 2026} Siga $P$ el punt de tall de les rectes $r_1:\left\{\begin{aligned}x&=1+t\\y&=0\\z&=-t\end{aligned}\right.$ \ i \ $r_2:\left\{\begin{aligned}x&=s\\y&=-1+s\\z&=0\end{aligned}\right.$ \ Calcula el punt $P'$ simètric de $P$ respecte del pla d'equació $2x+y-z=3$.

\ex{2} \pau{Madrid, ordinària 2026} Donats el pla $\pi:2x+2y-z=13$ i la recta \mbox{$r:\dfrac{x-2}{1}=\dfrac y1=\dfrac z4$}:\begin{tasks}(1)
  \task Troba el punt simètric del punt $P(2,\,0,\,0)$ respecte del pla $\pi$.
  \task Troba la distància entre el pla $\pi$ i la recta $r$.
\end{tasks}

\ex{2} \pau{Galícia, extraordinària 2026} Considera el punt $P(3,\,0,\,-1)$.\begin{tasks}(1)
  \task Calcula el punt $P'$ simètric de $P$ respecte del punt $Q(2,\,-1,\,2)$.
  \task Calcula el punt $P''$ simètric de $P$ respecte del pla $\pi:x+2y+z-8=0$.
\end{tasks}


% =====================================================
\seccio{10}{Àrees i volums}

\ex{1} Considera el triangle de vèrtexs $A(1,\,2,\,0)$, $B(3,\,0,\,1)$ i $C(0,\,1,\,2)$.\begin{tasks}(1)
  \task Calcula'n l'àrea.
  \task Calcula l'angle del vèrtex $A$.
  \task Calcula l'altura corresponent al costat $AB$.
\end{tasks}\recorda{$S=\tfrac12|\overrightarrow{AB}\times\overrightarrow{AC}|$. Per a l'angle d'un triangle no posem valor absolut al producte escalar.}

\ex{2} Calcula el volum del tetraedre de vèrtexs $A(1,\,0,\,0)$, $B(0,\,1,\,0)$, $C(0,\,0,\,1)$ i $D(2,\,2,\,2)$, i la seua altura des de $D$.

\ex{2} Troba $a$ perquè el tetraedre de vèrtexs $O$, $(1,\,0,\,0)$, $(0,\,1,\,0)$ i $(1,\,1,\,a)$ tinga volum 2.

\ex{2} Tres vèrtexs consecutius d'un paral·lelogram $ABCD$ són $A(1,\,0,\,1)$, $B(2,\,1,\,0)$ i $C(4,\,2,\,1)$. Troba el quart vèrtex i l'àrea del paral·lelogram.

\ex{2} El pla $\pi:2x+3y+6z=12$ talla els eixos en $A$, $B$ i $C$. Calcula el volum del tetraedre $OABC$, la distància de l'origen a $\pi$ i, amb aquestes dues dades, l'àrea del triangle $ABC$.

\ex{2} \pau{Cantàbria, extraordinària 2026} Considera els punts $A(0,\,5,\,-1)$ i $B(-3,\,11,\,-10)$ i la recta \mbox{$r:\dfrac{x-3}{1}=\dfrac{y+2}{-2}=\dfrac z3$}.\begin{tasks}(1)
  \task Estudia la posició relativa de $r$ i la recta que conté $A$ i $B$.
  \task Calcula l'àrea del triangle format pels vèrtexs $A$, $B$ i un punt genèric de $r$.
  \task Troba l'equació general o implícita del pla que conté $r$ i $A$.
\end{tasks}

\ex{2} \pau{Cantàbria, ordinària 2024} Siguen $A(6,\,2,\,-1)$, $B(3,\,0,\,5)$ i $C(-2,\,1,\,2)$ els vèrtexs d'un triangle.\begin{tasks}(1)
  \task Calcula els angles interns del triangle.
  \task Calcula l'àrea del triangle.
\end{tasks}

\ex{2} \pau{Extremadura, ordinària 2024} Donats els punts $A(1,\,2,\,1)$, $B(0,\,3,\,1)$ i $C(1,\,0,\,-1)$, determina:\begin{tasks}(1)
  \task un vector unitari i ortogonal als vectors $\overrightarrow{AB}$ i $\overrightarrow{AC}$;
  \task l'angle determinat per aquests vectors;
  \task l'àrea del triangle que formen $A$, $B$ i $C$.
\end{tasks}

\ex{2} \pau{Madrid, ordinària 2025} Siguen el punt $P(0,\,1,\,1)$ i el pla $\pi:x+y=2$.\begin{tasks}(1)
  \task Troba la distància del punt $P$ al pla $\pi$.
  \task Determina el punt $Q$ del pla $\pi$ la distància del qual a $P$ és igual que la distància de $P$ a $\pi$.
  \task Troba l'àrea del triangle format per $P$ i els punts de tall del pla $\pi$ amb els eixos de coordenades.
\end{tasks}

\ex{2} \pau{Navarra, extraordinària 2024} El punt $P(4,\,5,\,0)$ és el punt mig d'un costat d'un quadrat. El costat paral·lel a l'anterior està contingut en la recta $r:\left\{\begin{aligned}2x+2y+z&=0\\2x-z-2&=0\end{aligned}\right.$ \ Calcula els dos vèrtexs que determinen aquest segon costat.


% =====================================================
\seccio{11}{Llocs geomètrics}

\ex{1} Troba el pla mediador del segment d'extrems $A(2,\,1,\,-1)$ i $B(0,\,3,\,1)$ de dues maneres: imposant $d(X,A)=d(X,B)$ i amb el punt mig i el vector $\overrightarrow{AB}$.

\ex{2} Troba els plans bisectors dels plans $\pi_1:x+y-z=1$ i $\pi_2:x-y+z=3$.\recorda{Són els punts que equidisten dels dos plans: \mbox{$\dfrac{|A_1x+B_1y+C_1z+D_1|}{|\vec n_1|}=\dfrac{|A_2x+B_2y+C_2z+D_2|}{|\vec n_2|}$}.}

\ex{1} Troba el lloc geomètric dels punts que equidisten dels plans $2x-y+2z=1$ i $2x-y+2z=7$.

\ex{2} Escriu l'equació de l'esfera de centre $C(1,\,0,\,-2)$ que passa pel punt $P(3,\,2,\,-1)$ i la del pla tangent a l'esfera en $P$.

\ex{2} Troba l'esfera de centre $C(2,\,-1,\,1)$ tangent al pla $x+2y+2z+7=0$ i el punt de tangència.

\ex{3} Troba el lloc geomètric dels punts de l'espai que equidisten de $A(1,\,0,\,0)$, $B(0,\,1,\,0)$ i $C(0,\,0,\,1)$.

\ex{2} \pau{Andalusia, extraordinària 2024} Considera els punts $A(4,\,0,\,0)$ i $B(0,\,2,\,0)$. Calcula els punts del pla $OXZ$ que formen un triangle equilàter amb $A$ i $B$.

\ex{2} \pau{Andalusia, ordinària 2024} Considera les rectes $r:\left\{\begin{aligned}y&=0\\2x-z&=0\end{aligned}\right.$ \ i \ $s:\left\{\begin{aligned}x+y+7&=0\\z&=0\end{aligned}\right.$\begin{tasks}(1)
  \task Estudia la posició relativa de $r$ i $s$.
  \task Calcula l'equació del pla paral·lel a $r$ i a $s$ que equidista de totes dues rectes.
\end{tasks}

\ex{2} \pau{Castella i Lleó, ordinària 2026} Donats el pla $\pi_1:2x-3y+z=a$ i el pla $\pi_2$ determinat pel punt $P(0,\,2,\,4)$ i els vectors directors $\vec v_1=(0,\,2,\,6)$ i $\vec v_2=(1,\,0,\,b)$:\begin{tasks}(1)
  \task Calcula els valors de $a$ i $b$ perquè $\pi_1$ i $\pi_2$ siguen paral·lels i distints.
  \task Per a $a=1$ i $b=0$, determina les equacions paramètriques de la recta intersecció de $\pi_1$ i $\pi_2$.
  \task Per a $a=4$ i $b=-2$, determina els punts que estan a igual distància de $\pi_1$ i $\pi_2$.
\end{tasks}

\ex{2} \pau{Catalunya, ordinària 2025} Considera el pla $\pi:x+y=0$.\begin{tasks}(1)
  \task Calcula l'equació del pla $\pi'$ que és perpendicular a $\pi$ i conté els punts $P(1,\,-1,\,2)$ i $Q(3,\,-3,\,6)$.
  \task Calcula l'equació paramètrica de la recta continguda en $\pi'$ que conté els punts de $\pi'$ que estan a la mateixa distància de $P$ que de $Q$.
\end{tasks}

\ex{2} \pau{Catalunya, ordinària 2024} Considera els punts $A(1,\,2,\,3)$ i $B(-3,\,-2,\,3)$.\begin{tasks}(1)
  \task Calcula l'equació del pla $\pi$ que és perpendicular a la recta $AB$ i que passa pel punt mig entre $A$ i $B$. Justifica que aquest pla està format, precisament, pels punts $P$ que estan a igual distància de $A$ que de $B$.
  \task Calcula les distàncies de $A$ i de $B$ al pla $\pi$ i comprova que són iguals. És una casualitat?
  \task Siga $C(-7,\,6,\,3)$. El triangle $ABC$ és isòsceles? Calcula'n l'àrea.
\end{tasks}

\ex{2} \pau{Múrcia, extraordinària 2024} Considera els plans $x+y+z=-3$ i $x+y-z=3$ i la recta \mbox{$r:\dfrac{x-1}{2}=\dfrac{y+1}{1}=\dfrac z3$}.\begin{tasks}(1)
  \task Comprova que tots dos plans es tallen i calcula l'angle que formen.
  \task Estudia la posició relativa de la recta $r$ amb el pla $x+y-z=3$.
  \task Determina els punts de la recta $r$ que equidisten de tots dos plans.
\end{tasks}

\ex{2} \pau{Navarra, extraordinària 2025} Calcula l'equació del pla que equidista dels punts $A(3,\,3,\,5)$ i $B(1,\,-1,\,1)$. Calcula l'equació dels plans que disten 6 unitats del pla anterior.

\ex{2} \pau{Madrid, extraordinària 2024} Siguen els punts $P(1,\,-1,\,3)$ i $Q(2,\,1,\,-1)$.\begin{tasks}(1)
  \task Determina una equació del pla respecte del qual tots dos punts són simètrics.
  \task El segment $PQ$ és un dels tres costats d'un triangle en què la suma dels quadrats de les longituds dels costats és 34 i el tercer vèrtex es troba en la recta $r:x-2=y=z$. Calcula les coordenades del tercer vèrtex sabent que cap de les seues coordenades és nul·la.
\end{tasks}

\ex{2} \pau{País Basc, ordinària 2026} Siguen els punts $A(1,\,1,\,2)$ i $B(-1,\,0,\,3)$.\begin{tasks}(1)
  \task Troba l'equació del pla $\pi$ que compleix que $A$ i $B$ són simètrics respecte de $\pi$.
  \task Quina és la distància de $A$ a $\pi$?
\end{tasks}


% =====================================================
\seccio{12}{Problemes i qüestions}

\ex{2} Un raig de llum ix del punt $P(1,\,2,\,5)$ en la direcció del vector $(1,\,0,\,-1)$, es reflecteix en un espill situat en el pla $z=0$ i continua. Troba el punt on toca l'espill i l'equació del raig reflectit.\recorda{El raig reflectit és la recta que passa pel punt d'incidència i pel simètric de $P$ respecte de l'espill.}

\ex{3} Siguen $A(1,\,1,\,2)$ i $B(4,\,5,\,1)$, tots dos per damunt del pla $z=0$. Troba el punt $X$ del pla $z=0$ que fa mínima la suma $d(A,X)+d(X,B)$ i el valor d'aquesta suma mínima.

\ex{2} Raona si són certes o falses les afirmacions següents:\begin{tasks}(1)
  \task Dos plans perpendiculars a un mateix pla són sempre paral·lels entre si.
  \task Si una recta és perpendicular a un pla, és perpendicular a totes les rectes contingudes en el pla.
  \task La distància entre dues rectes que s'encreuen pot ser 0.
  \task Si $d(P,\pi)=0$, el punt $P$ és del pla $\pi$.
\end{tasks}

\ex{2} \pau{Astúries, ordinària 2026} Un mòbil es desplaça en línia recta des del punt $P(2,\,1,\,3)$ cap al punt $Q(0,\,-1,\,1)$. Quan arriba a la meitat de la distància de $P$ a $Q$ fa un gir i es desplaça per la recta que passa per $H(1,\,1,\,1)$.\begin{tasks}(1)
  \task Calcula les coordenades del punt en què fa el gir.
  \task Calcula la recta per la qual es desplaça el mòbil després del gir.
  \task Calcula l'angle de gir del mòbil.
\end{tasks}

\ex{2} \pau{Múrcia, extraordinària 2025} Un helicòpter situat en el punt $P(1,\,2,\,1)$ vol aterrar en el pla $\pi:x+y+3z=0$.\begin{tasks}(1)
  \task Calcula l'equació en forma contínua de la recta de la trajectòria que el porte al punt més pròxim de $\pi$.
  \task Troba aquest punt.
  \task Calcula la distància que ha de recórrer.
\end{tasks}

\ex{2} \pau{Astúries, ordinària 2025} S'estan construint dos ponts rectes en un tram d'autovia, un per a cada carril. Els ponts segueixen les equacions $r_1(t)=(2+t,\,-1-2t,\,3+2t)$ i $r_2(s)=(1+2s,\,4-s,\,4-2s)$.\begin{tasks}(1)
  \task Estudia si els ponts són paral·lels, es tallen o s'encreuen.
  \task L'empresa vol construir un pont de servei que els unisca i que siga el més curt possible. Quina longitud tindrà la via de servei? Indica'n els punts inicial i final.
\end{tasks}

\ex{2} \pau{Madrid, extraordinària 2026} Es consideren el punt $P(1,\,2,\,1)$ i el pla d'equació $2x-2y+z=1$. Es traça un cub amb una cara en el pla i un vèrtex en el punt $P$.\begin{tasks}(1)
  \task Determina la longitud del costat del cub i el seu volum.
  \task Troba un altre vèrtex del cub.
\end{tasks}

\ex{3} \pau{Aragó, extraordinària 2025} Es vol construir una estructura en forma de tetraedre la base de la qual té com a vèrtexs els punts $A(0,\,0,\,0)$, $B\left(2,\,0,\,\tfrac12\right)$ i $C\left(\tfrac32,\,3,\,1\right)$; el vèrtex superior, $D$, es troba en una biga recta entre els punts $E(0,\,1,\,3)$ i $F(3,\,2,\,3)$, és a dir, $D=E+\lambda\,\overrightarrow{EF}$ amb $\lambda\in[0,1]$.\begin{tasks}(1)
  \task Calcula el volum màxim d'aquesta estructura (totes les dades estan en metres).
  \task Tenint en compte que el volum d'una piràmide és un terç de l'àrea de la base per l'altura, calcula l'altura de l'estructura (des de $D$ fins a la base) si prenem $\lambda=1$.
\end{tasks}

\ex{2} \pau{Castella-la Manxa, ordinària 2026} En un parc públic s'està instal·lant un sistema de vigilància mitjançant sensors de moviment a banda i banda d'un mur recte, representat pel pla $\pi:x+y+z-6=0$. Un dels sensors està situat en el punt $P(1,\,2,\,6)$ i es vol instal·lar un segon sensor en el punt $P'$, simètric de l'anterior respecte del pla $\pi$.\begin{tasks}(1)
  \task Determina raonadament les coordenades del punt d'instal·lació del segon sensor, $P'$.
  \task Calcula la distància entre els dos sensors.
\end{tasks}

\ex{2} \pau{Castella-la Manxa, extraordinària 2025} Per a les festes del Corpus de Toledo s'instal·len tendals als carrers per on passa la processó. En un d'ells, els operaris col·loquen els punts de suport $A(0,\,1,\,-2)$, $B(1,\,2,\,0)$, $C(0,\,0,\,1)$ i $D(1,\,0,\,m)$, amb $m\in\mathbb R$.\begin{tasks}(1)
  \task Calcula el valor de $m$ perquè els quatre punts siguen coplanaris.
  \task Determina l'equació del pla $\pi$ que conté el tendal.
  \task Si els adorns florals han d'estar com a mínim a 1 metre de distància del tendal i s'ha col·locat un adorn de flors en el punt $P(1,\,2,\,3)$, estarà correctament situat?
\end{tasks}

\ex{2} \pau{Illes Balears, ordinària 2024} Considerem els punts $A(0,\,0,\,0)$, $B(2,\,-1,\,3)$ i $C(-1,\,2,\,1)$.\begin{tasks}(1)
  \task Calcula el punt $D$ tal que $ABDC$ siga un paral·lelogram.
  \task Calcula un dels punts $E$ de l'espai de manera que la recta $AE$ siga perpendicular al pla $ABC$ i la distància entre els punts $A$ i $E$ siga 1.
  \task Escriu l'equació d'un dels plans paral·lels al pla $ABC$ que disten una unitat d'aquest.
\end{tasks}

\vfill
\begin{center}\color{muted}\small\faCheckCircle[regular]\ \ Ànim! Comprova els teus resultats amb el laboratori mètric dels apunts interactius.\end{center}

\end{document}
